\chapter{Example Character Concepts}
\label{ch:example-concepts}

This chapter presents example character concepts to illustrate how the game's systems can create diverse and interesting heroes. These are \textbf{examples only}—not prescriptive templates or exhaustive lists. Use them for inspiration, as pre-generated characters, or as starting points for your own unique creations.

\section{Important Disclaimer}
\index{character concepts!disclaimer}

\textbf{These examples are provided for illustrative purposes only.} They demonstrate how the game's mechanics can support different character archetypes and play styles. You are encouraged to:
\begin{itemize}
\item Modify these concepts to fit your preferences
\item Create completely original characters
\item Mix and match elements from different examples
\item Work with your Game Master to develop unique concepts
\end{itemize}

The game system is designed to support a wide variety of character types beyond these examples.

\section{How to Use These Examples}
\index{character concepts!usage}

Each concept includes:
\begin{itemize}
\item \textbf{Concept Overview}: Narrative identity and role
\item \textbf{Traditional Analogues}: Common RPG classes (e.g., Fighter, Rogue, Wizard) that share similar themes — \textit{not} mechanical restrictions, just touchstones
\item \textbf{Mechanical Foundation}: Suggested starting capabilities
\item \textbf{Play Style}: How the character typically engages with challenges
\item \textbf{Development Path}: Potential growth directions
\item \textbf{Story Hooks}: Plot opportunities for the Game Master
\item \textbf{Build Blocks}: A \emph{30 XP} starting build, plus an optional \emph{34 XP} variant using Bonds/Complications (+4 XP)
\end{itemize}

\section{1. The Guardian (Fighter / Paladin / Knight)}
\index{character concepts!guardian}

\textbf{Concept}: A protector who stands between danger and those they've sworn to defend. \emph{Steel in hand, vow in heart.} Whether a chivalrous knight, a grizzled mercenary, or a holy warrior of the Everflame, the Guardian makes others safe through sacrifice.

\textbf{Traditional Analogues}: Fighter, Paladin, Barbarian (defensive), Knight, Samurai, Legionary

\textbf{Mechanical Foundation}:
\begin{itemize}
\item \textbf{Primary}: Body (strength, endurance), Spirit (willpower, resolve)
\item \textbf{Skills}: Melee, Athletics, Command, Endurance
\item \textbf{Talents}: Defensive Stance, Combat Reflexes, Conditioning
\end{itemize}

\textbf{Play Style}:
\begin{itemize}
\item Frontline combat and protection
\item Drawing attention away from allies (taunts, positioning)
\item Using presence and authority to control situations
\item Taking risks to protect others — absorbing Harm meant for teammates
\end{itemize}

\textbf{Development Path}:
\begin{itemize}
\item Increase defensive capabilities (armor, shields, parry)
\item Develop leadership skills (Command, Inspire)
\item Acquire better protective gear and magical wards
\item Learn area control abilities (shield wall, defensive line)
\end{itemize}

\textbf{Story Hooks}:
\begin{itemize}
\item Who or what are they protecting? A person, a place, an ideal?
\item What oath or duty drives them — sworn, inherited, or self‑imposed?
\item What happens if they fail in their protection?
\item What personal costs do they bear for their role (scars, lost loved ones, broken oaths)?
\end{itemize}

\paragraph{Build Blocks.}
\textbf{Starting Build (30 XP).}
\begin{itemize}
\item \textbf{Attributes} (Cost = rating $\times$ 3 XP): Body 3 (9), Spirit 2 (6), Wits 1 (3), Presence 1 (3) $\rightarrow$ \textbf{21 XP}
\item \textbf{Skills} (Cost = level $\times$ 2 XP): Melee 2 (4), Athletics 1 (2), Command 1 (2) $\rightarrow$ \textbf{8 XP}
\item \textbf{Total}: 29 XP (bank 1 XP)
\end{itemize}
\textbf{With Bonds/Complications (34 XP).}
\begin{itemize}
\item Add \textbf{Talent}: Combat Reflexes (2 XP) — +1 die on defense when surprised/flanked
\item Add \textbf{Talent}: Defensive Survival (3 XP) — +1 die to defense rolls while engaged in melee; once/scene convert first Harm 1 from melee to Fatigue
\item \textbf{Total}: 29 + 5 = 34 XP
\end{itemize}

\section{2. The Mage (Wizard / Sorcerer / Arcanist)}
\index{character concepts!mage}

\textbf{Concept}: A wielder of arcane power who reshapes reality through study, intuition, or pact. \emph{Candlesmoke, marginalia, and dangerous truths.} From the rational Wizard of Thepyrgos to the instinctive Sorcerer of the Bloodlands, magic flows through them — and sometimes through them.

\textbf{Traditional Analogues}: Wizard, Sorcerer, Warlock (pact magic), Arcanist, Runekeeper

\textbf{Mechanical Foundation}:
\begin{itemize}
\item \textbf{Primary}: Wits (theory, precision), Spirit (willpower, channeling)
\item \textbf{Skills}: Lore, Arcana, Investigation, (optionally) Sway (for social casting)
\item \textbf{Talents}: Spellcraft (freeform), Codex (runekeeper), Patron’s Symbol (invoker)
\end{itemize}

\textbf{Play Style}:
\begin{itemize}
\item Information gathering and magical analysis
\item Solving puzzles and mysteries with arcane insight
\item Using magic to gain advantages in combat and social scenes
\item Researching new spells and rituals between adventures
\item Managing Obligation (for Runekeepers) or Backlash (for free casters)
\end{itemize}

\textbf{Development Path}:
\begin{itemize}
\item Specialize in specific element or school of magic
\item Develop ritual casting or summoning abilities
\item Build a library, laboratory, or arcane network
\item Create unique spells or magical items
\end{itemize}

\textbf{Story Hooks}:
\begin{itemize}
\item What forbidden knowledge are they seeking?
\item How did they acquire their magical talent — bloodline, study, pact, accident?
\item What dangerous secrets might they uncover?
\item Who opposes their research (rival mages, witch‑hunters, the Church)?
\end{itemize}

\paragraph{Build Blocks.}
\textbf{Starting Build (30 XP).}
\begin{itemize}
\item \textbf{Attributes}: Wits 3 (9), Spirit 2 (6), Body 1 (3), Presence 1 (3) $\rightarrow$ \textbf{21 XP}
\item \textbf{Skills}: Lore 2 (4), Arcana 1 (2) $\rightarrow$ \textbf{6 XP}
\item \textbf{Total}: 27 XP (bank 3 XP)
\end{itemize}
\textbf{With Bonds/Complications (34 XP).}
\begin{itemize}
\item Add \textbf{Talent}: Spellcraft (6 XP) using banked 3 + 4 = 7 XP; bank 1 XP
\item \textbf{Total}: 27 + 6 = 33 XP (bank 1 XP)
\end{itemize}

\section{3. The Rogue (Thief / Assassin / Scout)}
\index{character concepts!rogue}

\textbf{Concept}: A shadow who moves where others cannot, taking what others cannot keep. \emph{Quiet footfalls, hawk eyes, and no loyalties but coin.} From the street‑wise cutpurse of Silkstrand to the deadly assassin of Zakov, the Rogue survives by wits, speed, and one hidden knife.

\textbf{Traditional Analogues}: Rogue, Thief, Assassin, Scout, Ranger (stealth focus), Spy

\textbf{Mechanical Foundation}:
\begin{itemize}
\item \textbf{Primary}: Wits (perception, quick thinking), Body (agility, stealth)
\item \textbf{Skills}: Stealth, Subterfuge, Investigation (or Streetwise), (optionally) Melee (light weapons)
\item \textbf{Talents}: Backstab, Elusive Dodge, Shadow Dance (later)
\end{itemize}

\textbf{Play Style}:
\begin{itemize}
\item Scouting ahead and gathering intelligence
\item Infiltrating secured areas and disabling traps
\item Ambush and skirmish tactics — striking from stealth
\item Finding paths, contacts, and secrets in cities and wilderness
\end{itemize}

\textbf{Development Path}:
\begin{itemize}
\item Improve stealth and backstab capabilities
\item Develop poison use or sabotage techniques
\item Build a network of informants and fences
\item Learn disguises, forgeries, or social infiltration
\end{itemize}

\textbf{Story Hooks}:
\begin{itemize}
\item What crime or secret forced them into the shadows?
\item What treasure or truth are they hunting?
\item How do they balance loyalty to allies with self‑preservation?
\item Who wants them dead — and why?
\end{itemize}

\paragraph{Build Blocks.}
\textbf{Starting Build (30 XP).}
\begin{itemize}
\item \textbf{Attributes}: Wits 3 (9), Body 2 (6), Presence 1 (3), Spirit 1 (3) $\rightarrow$ \textbf{21 XP}
\item \textbf{Skills}: Stealth 2 (4), Subterfuge 2 (4) $\rightarrow$ \textbf{8 XP}
\item \textbf{Total}: 29 XP (bank 1 XP)
\end{itemize}
\textbf{With Bonds/Complications (34 XP).}
\begin{itemize}
\item Add \textbf{Talent}: Backstab (6 XP) — needs 6 XP; bank 1 + 4 = 5 XP → not enough. Alternative: Add \textbf{Talent}: Elusive Dodge (6 XP) but same shortfall. Better: Add \textbf{Skill}: Perception 1 (2 XP) and \textbf{Talent}: Combat Reflexes (2 XP) and bank 2 XP. Or take a Minor Asset: Hidden Cache (4 XP) and bank 2 XP.
\item \textbf{Suggested}: Hidden Cache + bank 2 XP → total 33 XP.
\end{itemize}

\section{4. The Cleric (Priest / Healer / Exorcist)}
\index{character concepts!cleric}

\textbf{Concept}: A vessel of divine or spiritual power, channeling healing, truth, and judgment. \emph{Sacred flame, holy water, and a prayer on the lips.} From the Lampers of Ecktoria to the fog‑named exorcists of the Mistlands, the Cleric binds wounds and banishes darkness.

\textbf{Traditional Analogues}: Cleric, Priest, Paladin (healing focus), Exorcist, Shaman, Healer

\textbf{Mechanical Foundation}:
\begin{itemize}
\item \textbf{Primary}: Spirit (faith, willpower), Presence (authority, channeling)
\item \textbf{Skills}: Lore (theology, rites), Medicine, Sway (prayer, exhortation), (optionally) Melee (symbolic weapon)
\item \textbf{Talents}: Invoker (Patron’s Symbol) or Runekeeper (Thiasos + Codex of a divine Patron like Everflame, Oath of Flame & Light, Pale Shepherd)
\end{itemize}

\textbf{Play Style}:
\begin{itemize}
\item Healing and protecting allies
\item Purifying curses, diseases, and undead
\item Ley line and ritual magic (wards, consecrations)
\item Using authority of faith to command or persuade
\end{itemize}

\textbf{Development Path}:
\begin{itemize}
\item Deepen connection to Patron (more Rites, lower Obligation)
\item Develop exorcism and spirit‑banishing abilities
\item Build a temple, shrine, or hospice
\item Lead followers in religious/civic ceremonies
\end{itemize}

\textbf{Story Hooks}:
\begin{itemize}
\item What crisis of faith have they survived?
\item What heresy or supernatural threat haunts them?
\item How do they serve both the living and the dead?
\item What price do they pay for divine power?
\end{itemize}

\paragraph{Build Blocks.}
\textbf{Starting Build (30 XP).}
\begin{itemize}
\item \textbf{Attributes}: Spirit 3 (9), Presence 2 (6), Wits 1 (3), Body 1 (3) $\rightarrow$ \textbf{21 XP}
\item \textbf{Skills}: Medicine 2 (4), Lore 1 (2), Sway 1 (2) $\rightarrow$ \textbf{8 XP}
\item \textbf{Total}: 29 XP (bank 1 XP)
\end{itemize}
\textbf{With Bonds/Complications (34 XP).}
\begin{itemize}
\item Add \textbf{Talent}: Patron’s Symbol (4 XP) for a divine Patron (e.g., Oath of Flame & Light), using banked 1 + 4 = 5 XP; bank 1 XP
\item \textbf{Total}: 33 XP (bank 1 XP)
\end{itemize}

\section{5. The Bard (Diplomat / Skald / Courtier)}
\index{character concepts!bard}

\textbf{Concept}: A master of words, music, and social influence who changes minds and opens doors. \emph{A smile for the foyer, steel for the parlor.} From the Skald of Linn to the Courtier of Vhasia, the Bard wields charm, rumor, and reputation as weapons.

\textbf{Traditional Analogues}: Bard, Diplomat, Courtier, Skald, Spy (social), Negotiator

\textbf{Mechanical Foundation}:
\begin{itemize}
\item \textbf{Primary}: Presence (charm, performance), Wits (reading people, timing)
\item \textbf{Skills}: Sway, Performance, Insight, (optionally) Lore (customs, history)
\item \textbf{Talents}: Silver Tongue, Inspire, Network Builder (or Read Emotions)
\end{itemize}

\textbf{Play Style}:
\begin{itemize}
\item Social interaction and negotiation
\item Gathering information through rumors and contacts
\item Resolving conflicts without violence (or making violence fashionable)
\item Building alliances, manipulating factions
\end{itemize}

\textbf{Development Path}:
\begin{itemize}
\item Expand social network and influence
\item Develop economic or political power
\item Learn cultural specialties (heraldry, etiquette, languages)
\item Master inspirational or demoralizing performances
\end{itemize}

\textbf{Story Hooks}:
\begin{itemize}
\item What major conflict are they trying to resolve?
\item What secret do they know that could topple a dynasty?
\item How do they handle betrayal or failed negotiations?
\item What personal relationships affect their diplomacy?
\end{itemize}

\paragraph{Build Blocks.}
\textbf{Starting Build (30 XP).}
\begin{itemize}
\item \textbf{Attributes}: Presence 3 (9), Wits 2 (6), Spirit 1 (3), Body 1 (3) $\rightarrow$ \textbf{21 XP}
\item \textbf{Skills}: Sway 2 (4), Performance 1 (2), Insight 1 (2) $\rightarrow$ \textbf{8 XP}
\item \textbf{Total}: 29 XP (bank 1 XP)
\end{itemize}
\textbf{With Bonds/Complications (34 XP).}
\begin{itemize}
\item Add \textbf{Talent}: Inspire (3 XP) — once/scene, gain +1 die and potentially +1 Boon for bonded ally, using banked 1 + 4 = 5 XP; bank 2 XP
\item \textbf{Total}: 32 XP (bank 2 XP)
\end{itemize}

\section{6. The Ranger (Hunter / Tracker / Beastmaster)}
\index{character concepts!ranger}

\textbf{Concept}: A wilderness expert who navigates dangerous territories and hunts down quarry. \emph{Quiet footfalls, hawk eyes, and the long road.} From the Ykrul scout to the Valewood warden, the Ranger lives where roads end.

\textbf{Traditional Analogues}: Ranger, Hunter, Tracker, Beastmaster, Pathfinder, Wilderness Scout

\textbf{Mechanical Foundation}:
\begin{itemize}
\item \textbf{Primary}: Wits (perception, navigation), Body (agility, endurance)
\item \textbf{Skills}: Survival, Stealth, Nature (or Animal Handling), Archery or Melee
\item \textbf{Talents}: Wilderness Lore, Keen Senses, (optionally) Animal Companion
\end{itemize}

\textbf{Play Style}:
\begin{itemize}
\item Scouting ahead and gathering intelligence
\item Wilderness survival and navigation
\item Tracking creatures or people across terrain
\item Ambushing from concealment, using terrain
\item (Optional) Bonding with an animal companion
\end{itemize}

\textbf{Development Path}:
\begin{itemize}
\item Improve stealth and tracking abilities
\item Develop animal companion or nature magic
\item Master specific environments (forest, mountain, marsh)
\item Learn advanced survival techniques (poisoncraft, herbalism)
\end{itemize}

\textbf{Story Hooks}:
\begin{itemize}
\item What uncharted territory are they exploring?
\item What secrets have they discovered in the wild?
\item How do they balance civilization and wilderness?
\item What prey (beast or man) still eludes them?
\end{itemize}

\paragraph{Build Blocks.}
\textbf{Starting Build (30 XP).}
\begin{itemize}
\item \textbf{Attributes}: Wits 3 (9), Body 2 (6), Spirit 1 (3), Presence 1 (3) $\rightarrow$ \textbf{21 XP}
\item \textbf{Skills}: Survival 2 (4), Stealth 2 (4) $\rightarrow$ \textbf{8 XP}
\item \textbf{Total}: 29 XP (bank 1 XP)
\end{itemize}
\textbf{With Bonds/Complications (34 XP).}
\begin{itemize}
\item Add \textbf{Talent}: Keen Senses (2 XP) — +1 die to perception
\item Add \textbf{Talent}: Wilderness Lore (2 XP) — auto‑find food/water in hospitable biomes
\item \textbf{Total}: 29 + 4 = 33 XP (bank 1 XP)
\end{itemize}

\section{7. The Survivor (Veteran / Outcast / Hardened)}
\index{character concepts!survivor}

\textbf{Concept}: Someone who has endured hardship and developed resilience. \emph{Scars are maps; read them well.} From the war‑torn mercenary to the plague survivor, the Survivor won’t break — they’ve already been broken and rebuilt.

\textbf{Traditional Analogues}: Barbarian (resilience focus), Veteran, Ranger (endurance), Drifter, Refugee

\textbf{Mechanical Foundation}:
\begin{itemize}
\item \textbf{Primary}: Spirit (willpower, resolve), Body (toughness)
\item \textbf{Skills}: Endurance, Survival, (optionally) Perception (danger sense), Medicine (field)
\item \textbf{Talents}: Iron Stomach, Conditioning, Second Wind
\end{itemize}

\textbf{Play Style}:
\begin{itemize}
\item Enduring difficult conditions (lack of food, exposure, poison)
\item Overcoming physical and mental challenges through grit
\item Using experience to avoid dangers
\item Helping others survive hardships
\item Taking hits that would fell others and still standing
\end{itemize}

\textbf{Development Path}:
\begin{itemize}
\item Improve physical and mental resilience (higher Fatigue limit, Harm reduction)
\item Develop survival‑related skills (traps, scavenging)
\item Acquire better equipment and resources
\item Learn to teach survival to others (leadership)
\end{itemize}

\textbf{Story Hooks}:
\begin{itemize}
\item What trauma or hardship have they survived?
\item How has their past shaped their present?
\item What survival skills have saved them repeatedly?
\item How do they help others facing similar challenges?
\end{itemize}

\paragraph{Build Blocks.}
\textbf{Starting Build (30 XP).}
\begin{itemize}
\item \textbf{Attributes}: Spirit 3 (9), Body 2 (6), Wits 1 (3), Presence 1 (3) $\rightarrow$ \textbf{21 XP}
\item \textbf{Skills}: Endurance 2 (4), Survival 2 (4) $\rightarrow$ \textbf{8 XP}
\item \textbf{Total}: 29 XP (bank 1 XP)
\end{itemize}
\textbf{With Bonds/Complications (34 XP).}
\begin{itemize}
\item Add \textbf{Talent}: Iron Stomach (3 XP) — immunity to mundane poisons and spoiled food; halve complications from toxic sources
\item \textbf{Total}: 29 + 3 = 32 XP (bank 2 XP)
\end{itemize}

\section{8. The Innovator (Artificer / Engineer / Tinker)}
\index{character concepts!innovator}

\textbf{Concept}: A creative problem-solver who finds new solutions through craft and invention. \emph{Blueprints on napkins, tomorrow in your pocket.} From the Aeler engineer to the Gnome clockmaker, the Innovator builds what others only dream.

\textbf{Traditional Analogues}: Artificer, Engineer, Tinker, Alchemist, Inventor, Savant

\textbf{Mechanical Foundation}:
\begin{itemize}
\item \textbf{Primary}: Wits (design, analysis), Body (fine manipulation)
\item \textbf{Skills}: Craft, Tinker, Mechanics, (optionally) Alchemy, Investigation
\item \textbf{Talents}: Technical Expert, Quick Study, (optionally) Magical Laboratory asset
\end{itemize}

\textbf{Play Style}:
\begin{itemize}
\item Solving problems with unique inventions
\item Creating or repairing specialized items
\item Providing services others cannot (lockpicking, trap disarm, device creation)
\item Using knowledge to gain advantages in planning
\end{itemize}

\textbf{Development Path}:
\begin{itemize}
\item Master their specialty area (weapons, armor, alchemy, constructs)
\item Develop related capabilities (explosives, poisons, medical devices)
\item Build reputation and clientele (workshop asset)
\item Create unique inventions or discoveries (blueprints)
\end{itemize}

\textbf{Story Hooks}:
\begin{itemize}
\item What makes their specialty unique or valuable?
\item How did they acquire their special skills (guild, apprenticeship, self‑taught)?
\item What problems require their specific expertise?
\item Who seeks to control or exploit their abilities?
\end{itemize}

\paragraph{Build Blocks.}
\textbf{Starting Build (30 XP).}
\begin{itemize}
\item \textbf{Attributes}: Wits 3 (9), Body 2 (6), Presence 1 (3), Spirit 1 (3) $\rightarrow$ \textbf{21 XP}
\item \textbf{Skills}: Craft 2 (4), Tinker 2 (4) $\rightarrow$ \textbf{8 XP}
\item \textbf{Total}: 29 XP (bank 1 XP)
\end{itemize}
\textbf{With Bonds/Complications (34 XP).}
\begin{itemize}
\item Add \textbf{Talent}: Quick Study (3 XP) — +1 die to learn or adapt from observation, using banked 1 + 4 = 5 XP; bank 2 XP
\item \textbf{Total}: 32 XP (bank 2 XP)
\end{itemize}

\section{9. The Summoner (Pact‑Whisperer / Necromancer / Shaman)}
\index{character concepts!summoner}

\textbf{Concept}: A binder of spirits, elementals, and otherworldly entities. \emph{Circles of salt, whispered names, and bargains struck in shadows.} From the necromancer of the Mistlands to the shaman of the Steppe, the Summoner commands powers that others fear.

\textbf{Traditional Analogues}: Necromancer, Conjurer, Summoner, Shaman, Spirit Caller, Demonologist

\textbf{Mechanical Foundation}:
\begin{itemize}
\item \textbf{Primary}: Spirit (willpower, binding), Wits (rituals, names)
\item \textbf{Skills}: Arcana, Lore (spirits), (optionally) Command, Survival
\item \textbf{Talents}: Pact‑Whisperer (Lesser/Greater), Spirit Link, (optionally) Codex of Names
\end{itemize}

\textbf{Play Style}:
\begin{itemize}
\item Summoning creatures to fight, scout, or work
\item Binding entities to solve problems (guard, research, transport)
\item Negotiating with spirits and outsiders (social challenges)
\item Managing Leash clocks and Boon Finesse
\end{itemize}

\textbf{Development Path}:
\begin{itemize}
\item Unlock higher Cap summons (Lesser → Greater)
\item Develop specializations (combat, scout, utility)
\item Build a Spirit Bond Clock with favored entities
\item Learn true names for powerful outsiders (Prestige talent)
\end{itemize}

\textbf{Story Hooks}:
\begin{itemize}
\item What entity are they bound to or seeking?
\item What price did they pay for their first summoning?
\item How do others react to their “dark” magic?
\item Who hunts them for their forbidden practices?
\end{itemize}

\paragraph{Build Blocks.}
\textbf{Starting Build (30 XP).}
\begin{itemize}
\item \textbf{Attributes}: Spirit 3 (9), Wits 2 (6), Body 1 (3), Presence 1 (3) $\rightarrow$ \textbf{21 XP}
\item \textbf{Skills}: Arcana 2 (4), Lore 1 (2) $\rightarrow$ \textbf{6 XP}
\item \textbf{Total}: 27 XP (bank 3 XP)
\end{itemize}
\textbf{With Bonds/Complications (34 XP).}
\begin{itemize}
\item Add \textbf{Talent}: Pact‑Whisperer (2 XP) — access to summoning
\item Add \textbf{Talent}: Lesser Pactwright (2 XP) — can summon Cap 1 spirits
\item \textbf{Total}: 27 + 4 = 31 XP (bank 3 XP; can add Familiar (2 XP) later or upgrade to Greater Pactwright)
\end{itemize}

\section{10. The Commander (Warlord / Tactician / Leader)}
\index{character concepts!commander}

\textbf{Concept}: A leader who directs allies, controls battlefields, and shapes conflicts through strategy. \emph{Orders in a whisper, victory in formation.} From the Condotta captain to the Ekstorian legate, the Commander wins before blades are drawn.

\textbf{Traditional Analogues}: Warlord, Marshal, Tactician, Commander, Strategist, Banneret

\textbf{Mechanical Foundation}:
\begin{itemize}
\item \textbf{Primary}: Presence (authority, command), Wits (tactics, observation)
\item \textbf{Skills}: Command, Tactics, (optionally) Sway (morale), Investigation (battlefield analysis)
\item \textbf{Talents}: Inspire, Command Presence, Tactical Movement (for allies)
\end{itemize}

\textbf{Play Style}:
\begin{itemize}
\item Directing allies in combat (bonus dice, improved Position)
\item Planning and executing tactical maneuvers
\item Leading followers and assets in mass action
\item Using authority to intimidate or persuade
\end{itemize}

\textbf{Development Path}:
\begin{itemize}
\item Increase follower limit and effectiveness
\item Develop area‑buff abilities (morale, defensive formation)
\item Acquire strategic assets (spy network, war chest)
\item Unlock Prestige talents like Battlefield Mastery
\end{itemize}

\textbf{Story Hooks}:
\begin{itemize}
\item What army or company did they lead before?
\item What strategic defeat haunts them?
\item Who are their rivals in the art of war?
\item What cause do they fight for — coin, country, or creed?
\end{itemize}

\paragraph{Build Blocks.}
\textbf{Starting Build (30 XP).}
\begin{itemize}
\item \textbf{Attributes}: Presence 3 (9), Wits 2 (6), Body 1 (3), Spirit 1 (3) $\rightarrow$ \textbf{21 XP}
\item \textbf{Skills}: Command 2 (4), Tactics 2 (4) $\rightarrow$ \textbf{8 XP}
\item \textbf{Total}: 29 XP (bank 1 XP)
\end{itemize}
\textbf{With Bonds/Complications (34 XP).}
\begin{itemize}
\item Add \textbf{Talent}: Command Presence (4 XP) — +1 die on leadership/intimidation rolls, using banked 1 + 4 = 5 XP; bank 1 XP
\item \textbf{Total}: 33 XP (bank 1 XP)
\end{itemize}

\begin{tcolorbox}[colback=gray!5!white,colframe=gray!70!black,title=Reskinning the Paths: Runekeeper and Invoker as Familiar Concepts]
    The magic systems in \textit{Fate's Edge} are designed to be reimagined. Neither \textbf{Runekeeper} (Thiasos + Codex) nor \textbf{Invoker} (Patron’s Symbol) is locked into a single narrative identity. Below are common reskins — use the same mechanics, change the fiction.
    
    \subsection*{Runekeeper – The Devoted Agent}
    A Runekeeper has sworn themselves to a single Patron, learning its Rites and bearing its Obligation. This framework can represent:
    
    \begin{itemize}
      \item \textbf{Paladin / Crusader}: Sworn to the \textit{Oath of Flame \& Light} or the \textit{Inquisitor Prime}. Rites become prayers, smites, and consecrations. The Codex is a prayer book or a vow‑etched gauntlet.
      \item \textbf{Warlock / Pact‑Bearer}: Bound to a distant, alien, or dangerous Patron (e.g., Khemesh, Malachai, the Ninth). Rites are granted, not learned. Obligation is the steady pull of a bargain.
      \item \textbf{Cleric / Priest}: Devoted to the Everflame, the Pale Shepherd, or other divine powers. Rites are liturgies. The Patron’s Gift is a blessing or a relic.
      \item \textbf{Dark Knight / Death Knight}: A fallen or morally ambiguous warrior who draws power from a grim Patron (Carrion‑King, Mor’iraath). Melee and Rites mix; Obligation is a slow doom.
    \end{itemize}
    
    \noindent \textit{Mechanical core unchanged:} 1 action Rites, Obligation track, Patron’s Gift imbuement, Push It option.
    
    \subsection*{Invoker – The Professional Occultist}
    An Invoker does not fully commit to a single Patron. Instead, they carry \textbf{Symbols} (often tools, badges, or relics) of several Powers and perform rituals to invoke their Rites. This framework can represent:
    
    \begin{itemize}
      \item \textbf{Exorcist / Witch Hunter}: Carries symbols of several protective Patrons (The Sealed Gate, Oath of Flame, Gallow’s Bell). Rites are wards, banishments, and purification rituals. \emph{Crack the Seal} is an emergency exorcism.
      \item \textbf{Alienist / Scholar of the Unseen}: Studies forbidden Patrons (The Ninth, Khemesh, Mor’iraath) from a safe distance — “safe” being relative. Uses Symbol libraries to invoke their power without full devotion. Rituals are experimental and dangerous.
      \item \textbf{Consulting Occultist / Paranormal Investigator}: Carries a briefcase or satchel of Symbols — a bell from the Mistlands, a sun‑seal from Ecktoria, a hag’s stone from Aelinnel. Hired to solve magical problems. Rituals are filed like work orders; \emph{Crack the Seal} is “burning a favor.”
      \item \textbf{Ritualist / Ceremonial Mage}: Believes that power lies in procedure, not faith. Each Symbol is a “key” — a geometric diagram, a calendar, a consecrated blade. Rites are slow, deliberate, and precise. The Invoker is a technician, not a worshipper.
    \end{itemize}
    
    \noindent \textit{Mechanical core unchanged:} Ritual time, Obligation per invocation, Symbol display reduces cost, \emph{Crack the Seal} for instant effect.
    
    \subsection*{Mixing and Blurring Lines}
    Characters may also blend these frameworks — a Runekeeper who occasionally carries a second Symbol for emergency rites, or an Invoker who one day pledges fully to a Patron and becomes a Runekeeper. The system does not penalize narrative growth; simply retrain Talents during downtime (\S\ref{ch:talents}) when the fiction justifies it.
    
    \noindent \textit{Remember:} The fiction leads. If you want to play a Paladin, describe your Runekeeper’s Rites as prayers. If you want a Vodouisant or Hoodoo practitioner, reflavor Patron’s Symbols as sacred bundles, veves, or spirit vessels. If you want a monster hunter, carry a satchel of monster‑specific Symbols (silver stake, iron bell, cold iron nail). The mechanics stay the same; the story breathes.
    \end{tcolorbox}
    
\section{Creating Your Own Concept}
\index{character concepts!creation}

\subsection*{Start with Narrative}
\begin{itemize}
\item What is your character's background and motivation?
\item What role do they play in their community or society?
\item What relationships are important to them?
\item What goals are they pursuing?
\end{itemize}

\subsection*{Add Mechanical Support}
\begin{itemize}
\item Choose attributes that support your concept
\item Select skills that reflect their training and experience
\item Consider talents that provide unique capabilities
\item Think about assets that represent their resources
\end{itemize}

\subsection*{Consider Group Role}
\begin{itemize}
\item How does your concept complement other party members?
\item What gaps in group capability can you fill?
\item What unique contributions can you make?
\item How will you work with other characters?
\end{itemize}

\subsection*{Plan for Growth}
\begin{itemize}
\item What short-term improvements make sense?
\item What long-term development aligns with your concept?
\item How might your character change over time?
\item What legacy do you want to build?
\end{itemize}

\begin{tcolorbox}[colback=green!5!white,colframe=green!75!black,title=Character Concept Worksheet,fonttitle=\bfseries]
\textbf{Narrative Elements:}
\begin{itemize}
\item Concept: \_\_\_\_\_\_\_\_\_\_\_\_\_\_\_\_\_\_\_\_\_\_
\item Traditional Analogues: \_\_\_\_\_\_\_\_\_\_\_\_\_\_\_\_
\item Motivation: \_\_\_\_\_\_\_\_\_\_\_\_\_\_\_\_\_\_\_\_\_\_
\item Background: \_\_\_\_\_\_\_\_\_\_\_\_\_\_\_\_\_\_\_\_\_\_
\item Relationships: \_\_\_\_\_\_\_\_\_\_\_\_\_\_\_\_\_\_\_\_\_\_
\end{itemize}

\textbf{Mechanical Foundation:}
\begin{itemize}
\item Primary Attributes: \_\_\_\_\_\_\_\_\_\_\_\_\_\_\_\_\_\_\_\_\_\_
\item Key Skills: \_\_\_\_\_\_\_\_\_\_\_\_\_\_\_\_\_\_\_\_\_\_
\item Starting Talents: \_\_\_\_\_\_\_\_\_\_\_\_\_\_\_\_\_\_\_\_\_\_
\item Initial Assets: \_\_\_\_\_\_\_\_\_\_\_\_\_\_\_\_\_\_\_\_\_\_
\end{itemize}

\textbf{Development Plan:}
\begin{itemize}
\item Short-term goals: \_\_\_\_\_\_\_\_\_\_\_\_\_\_\_\_\_\_\_\_\_\_
\item Long-term vision: \_\_\_\_\_\_\_\_\_\_\_\_\_\_\_\_\_\_\_\_\_\_
\end{itemize}
\end{tcolorbox}

\section{Final Notes}

Remember that these examples are starting points, not limitations. The most interesting characters often combine elements from multiple concepts or create entirely new approaches. Work with your Game Master to ensure your character concept fits the campaign and provides engaging story opportunities.

The best characters are those that you find interesting to play and that contribute to an enjoyable experience for everyone at the table.
